%==================================================================
% Akif Faruk Kazi: Resume (industry, one page), "void" edition
% Compile with pdfLaTeX (Overleaf default). Run twice if the page
% frame does not appear on the first compile (Overleaf does this).
%
% Look: sterile, cool, near-empty UI. White page with a faint ice
% glow, one navy "selected" tab, thin letter-spaced Montserrat for
% display type, hairline rules with square markers, viewfinder
% brackets in the page corners.
%
% ATS notes (unchanged from resume.tex)
%   - Single column. Every piece of text is ordinary text written
%     in reading order; the tab, rules, squares, glow and brackets
%     are vector graphics with no text in them.
%   - Standard section names: Experience, Education, Skills.
%   - T1 fonts + glyphtounicode, ligatures and hyphenation off.
%   - Links written out in full.
%==================================================================

\documentclass[a4paper,10pt]{article}

%----------PAGE----------
\usepackage[
    left=1.15cm,
    right=1.15cm,
    top=1.05cm,
    bottom=0.9cm
]{geometry}

%----------TYPE----------
% Montserrat (display): geometric, wide, close to the thin UI type
% of sci-fi menus. Source Sans Pro (body): open and highly legible.
\usepackage[T1]{fontenc}
\usepackage[utf8]{inputenc}
\usepackage[light,semibold]{montserrat}
\usepackage[default,semibold,scale=1.08]{sourcesanspro}

% T1 is pinned so the display face resolves the same way in every
% engine (pdfLaTeX on Overleaf, XeTeX for local previews).
\newcommand{\display}{\fontencoding{T1}\fontfamily{Montserrat-TLF}\selectfont}
\newcommand{\displaylight}{\display\fontseries{l}\selectfont}
\newcommand{\displaysemi}{\display\fontseries{sb}\selectfont}

\usepackage{iftex}
\ifPDFTeX
    \input{glyphtounicode}
    \pdfgentounicode=1
\fi

\usepackage{microtype}
\ifPDFTeX
    \DisableLigatures{encoding = *, family = *}
    \newcommand{\track}[2]{\textls[#1]{#2}}
\else
    \newcommand{\track}[2]{#2}
\fi

\usepackage{xcolor}
\usepackage{enumitem}
\usepackage{titlesec}
\usepackage{tikz}
\usepackage{eso-pic}
\usepackage[hidelinks]{hyperref}

%----------PDF METADATA----------
\hypersetup{
    pdftitle={Akif Faruk Kazi - Resume},
    pdfauthor={Akif Faruk Kazi},
    pdfsubject={Interaction Designer, XR Designer},
    pdfkeywords={Interaction Design, UX Design, XR Design, Spatial
        Interaction, Prototyping, Usability Testing, Design Systems,
        Figma, Unity, Blender, WebGL}
}

%----------COLOUR----------
% Cold, desaturated slate with one deep navy accent.
\definecolor{navy}{HTML}{174A70}    % selected tab, markers
\definecolor{ink}{HTML}{1B2631}     % name, company names
\definecolor{body}{HTML}{3A4653}    % body text
\definecolor{muted}{HTML}{6F7E8C}   % dates, locations
\definecolor{rule}{HTML}{C9D4DD}    % hairlines
\definecolor{ice}{HTML}{E4EDF4}     % background glow
\definecolor{frost}{HTML}{B9C9D6}   % tab extension, brackets

%----------SPACING----------
% Tighten these first if a new entry pushes the page over.
\newlength{\entrygap}\setlength{\entrygap}{9pt}
\newlength{\bulletgap}\setlength{\bulletgap}{3pt}
\linespread{1.12}

\pagestyle{empty}
\setlength{\parindent}{0pt}
\setlength{\parskip}{0pt}
\raggedright
\hyphenpenalty=10000
\exhyphenpenalty=10000
\urlstyle{same}
\color{body}

%----------PAGE BACKGROUND----------
% A faint cold glow from the top right and viewfinder brackets in
% the four corners. Graphics only; nothing here is text.
\AddToShipoutPictureBG{%
    \begin{tikzpicture}[remember picture, overlay]
        % Soft glow built from 30 flat rings (no shading objects, so
        % every PDF viewer, printer and parser handles it).
        \foreach \i in {0,...,29} {
            \pgfmathsetmacro{\pct}{100*(\i+1)/30}
            \fill[ice!\pct!white]
                (current page.north east) circle ({13cm - \i*0.42cm});
        }
        \draw[frost, line width=0.7pt]
            ([xshift=6mm, yshift=-12mm]current page.north west) --
            ([xshift=6mm, yshift=-6mm]current page.north west) --
            ([xshift=12mm, yshift=-6mm]current page.north west);
        \draw[frost, line width=0.7pt]
            ([xshift=-12mm, yshift=-6mm]current page.north east) --
            ([xshift=-6mm, yshift=-6mm]current page.north east) --
            ([xshift=-6mm, yshift=-12mm]current page.north east);
        \draw[frost, line width=0.7pt]
            ([xshift=6mm, yshift=12mm]current page.south west) --
            ([xshift=6mm, yshift=6mm]current page.south west) --
            ([xshift=12mm, yshift=6mm]current page.south west);
        \draw[frost, line width=0.7pt]
            ([xshift=-12mm, yshift=6mm]current page.south east) --
            ([xshift=-6mm, yshift=6mm]current page.south east) --
            ([xshift=-6mm, yshift=12mm]current page.south east);
    \end{tikzpicture}%
}

%----------MARKERS----------
\newcommand{\sq}[1][2.2pt]{\raisebox{0.32ex}{\color{navy}\rule{#1}{#1}}}
\newcommand{\sep}{\unskip\enspace{\color{rule}\textbar}\enspace\ignorespaces}

%----------SELECTED TAB----------
% Navy block with a slanted edge running into a pale extension,
% the "active menu item" of the page. The label is ordinary text.
\newcommand{\selectedtab}[1]{%
    \begin{tikzpicture}[baseline=(t.base)]
        % Invisible copy to measure the box; draws no text.
        \node[inner xsep=8pt, inner ysep=4.5pt, anchor=west]
            (t) at (0,0) {\phantom{\displaysemi\small\track{140}{#1}}};
        \fill[frost!50] (t.south east) rectangle ([xshift=2.8cm]t.north east);
        \fill[navy] (t.south west) -- (t.north west) --
            ([xshift=9pt]t.north east) -- (t.south east) -- cycle;
        % The one real copy of the label.
        \node[inner xsep=8pt, inner ysep=4.5pt, anchor=west, text=white]
            at (0,0) {\displaysemi\small\track{140}{#1}};
    \end{tikzpicture}%
}

%----------SECTIONS----------
% Unselected-menu style: letter-spaced Montserrat, then a hairline
% that starts and ends on a small square.
\newcommand{\sectionhead}[1]{%
    \makebox[\linewidth][l]{%
        {\color{ink}\displaysemi\small\track{160}{\MakeUppercase{#1}}}%
        \hspace{8pt}\sq[2pt]\hspace{1pt}%
        {\color{rule}\leaders\hbox{\rule[0.62ex]{1pt}{0.5pt}}\hfill}%
        \hspace{1pt}\sq[2pt]}%
}
\titleformat{\section}{}{}{0pt}{\sectionhead}
\titlespacing*{\section}{0pt}{12pt}{6pt}

%----------LISTS----------
\setlist[itemize]{
    leftmargin=1.1em,
    labelsep=0.55em,
    itemsep=\bulletgap,
    topsep=4pt,
    parsep=0pt,
    partopsep=0pt,
    label={\sq[2.4pt]}
}

%----------ENTRY----------
% \entry{Organisation}{Dates}{Role}{Location}
% Line 1: organisation, dates. Line 2: role (navy), location.
\newcommand{\entry}[4]{%
    {\color{ink}\textbf{#1}}\hfill{\color{muted}#2}\par
    {\color{navy}#3}\hfill{\color{muted}#4}\par
}

%==================================================================
\begin{document}

%----------HEADER----------
{\color{ink}\fontsize{27}{31}\selectfont\displaylight\track{110}{AKIF FARUK KAZI}}\par
\vspace{6pt}
\selectedtab{INTERACTION DESIGNER \enspace|\enspace XR DESIGNER}\par
\vspace{6pt}
{\small\color{body}Mumbai, India
\sep +91 98677 35452
\sep \href{mailto:akif.kazi@icloud.com}{akif.kazi@icloud.com}
\sep \href{https://www.linkedin.com/in/akif-faruk-kazi/}{linkedin.com/in/akif-faruk-kazi}
\sep \href{https://akifkazi.notion.site/portfolio}{akifkazi.notion.site/portfolio}}\par

%----------EXPERIENCE----------
\section{Experience}

\entry{Rosatom}{May 2026 \textendash{} Present}{Design Intern}{Moscow, Russia (Remote)}
\begin{itemize}
    \item Designed a browser-based digital twin for monitoring industrial equipment on the go, beyond the control room.
    \item Developed a sci-fi inspired HUD interface with responsive, orientation-aware layouts from desktop to mobile.
    \item Architected a modular, configuration-driven platform so new machines map to their sensor tags without redesign.
    \item Shipped WebGL builds that load in under 7 seconds with per-texture compression and IL2CPP code stripping.
\end{itemize}
\vspace{\entrygap}

\entry{Indian Institute of Technology Jodhpur}{July 2025 \textendash{} Present}{Teaching Assistant, Engineering Design}{Jodhpur, India}
\begin{itemize}
    \item Mentor and give design feedback to 48 students from 13 engineering disciplines building products in teams.
\end{itemize}
\vspace{\entrygap}

\entry{Siemens}{January 2025 \textendash{} July 2025}{Design Intern, FT UX Advanced Experience Research}{Bengaluru, India}
\begin{itemize}
    \item Designed scalable UI components for the COMOS Walkinside digital twin, showcased at UX@Siemens Germany.
    \item Contributed reusable XR interaction patterns proposed for Siemens iX, the company's open-source design system.
    \item Researched industrial XR heads-up displays and evaluated concepts with stakeholders, refining information hierarchy, placement, interaction behaviour and component visibility.
    \item Helped build a controlled test environment for a patent-pending VR interaction technique.
\end{itemize}
\vspace{\entrygap}

\entry{Elevar}{May 2023 \textendash{} September 2023}{UI/UX Design Intern}{Mumbai, India}
\begin{itemize}
    \item Built a token-based Figma design system (colour, spacing, radius) with responsive component variants, and used it to redesign and code a responsive Contact page in HTML/CSS.
    \item Analysed user behaviour in Google Analytics and Microsoft Clarity to prioritise usability fixes.
\end{itemize}

%----------EDUCATION----------
\section{Education}

\entry{Indian Institute of Technology Jodhpur}{2025 \textendash{} Present}{Master of Design, XR Design \sep CGPA 7.88}{Jodhpur, India}
\vspace{\entrygap}

\entry{Somaiya Vidyavihar University}{2021 \textendash{} 2025}{Bachelor of Design, Interaction Design \sep CGPA 9.23}{Mumbai, India}

%----------SKILLS----------
\section{Skills}

\textbf{\color{ink}Design:} Interaction Design, XR Design, UX Research, Usability Testing, Prototyping, Design Systems, Data Visualisation\par
\vspace{4pt}
\textbf{\color{ink}Tools:} Figma, Unity, Blender, Adobe Creative Suite, Affinity, WebGL, HTML/CSS, Google Analytics, Microsoft Clarity\par

%----------ACHIEVEMENTS----------
\section{Achievements}

AIR 57, NID GD 2025 \sep AIR 417, CEED 2025 \sep AIR 64, UCEED 2021\par

%----------LEADERSHIP----------
\section{Leadership and Activities}

\entry{TEDxIIMLucknow}{September 2024}{Design Team Member}{IIM Lucknow (Remote)}
\vspace{\entrygap}

\entry{Nazariya '23 Design Fest}{October 2023 \textendash{} November 2023}{Graphic Design Head}{Somaiya Vidyavihar University}
\vspace{\entrygap}

\entry{Ycenter}{September 2021 \textendash{} November 2021}{Social Innovation Fellow}{Philadelphia, USA (Remote)}

\end{document}
